%!TeX root=MemoriaTFG.tex

\chapter{Prototipo DermApIxel}
\label{cap:anexoh}

Este anexo describe el prototipo \emph{DermApIxel}, sistema
integrado de apoyo al diagnóstico dermatológico construido como
entorno operativo de los módulos derivados del trabajo
experimental. La versión inicial documentada en las
secciones~\ref{sec:dermapixel-motivacion}--\ref{sec:dermapixel-discusion}
opera con seis módulos de inteligencia artificial
(M1--M6), estado correspondiente al despliegue \emph{end-to-end}
de 2026-04-05. La evolución posterior incorpora tres módulos
castellanos adicionales (M7, M9, M10) más un módulo de
recuperación sobre DermapixelAI~1.0 (M4-bis) y un ensemble
ponderado (M11), descritos en la
sección~\ref{sec:dermapixel-m11}. El prototipo no constituye
objeto experimental del estudio principal del documento
(sección~\ref{sec:alcance}), pero materializa la integración
técnica entre los modelos evaluados en el
capítulo~\ref{cap:resultados} y un sistema operativo accesible
desde navegador.

\section{Motivación arquitectónica}
\label{sec:dermapixel-motivacion}

La traslación de los hallazgos cuantitativos del benchmark a un
entorno operativo plantea requisitos no satisfechos por las
herramientas estándar de evaluación (notebooks de investigación,
scripts CLI). Cuatro requisitos vertebran el diseño del prototipo:
(i)~separación física entre el equipo cliente (PC del clínico, sin
GPU) y el servidor de inferencia (DGX Spark, con GPU);
(ii)~comunicación asíncrona entre ambos para evitar bloqueos en la
interfaz durante inferencias de varios segundos; (iii)~persistencia
estructurada de imagen original, máscara, conceptos, casos
similares, razonamiento textual y validación posterior por
dermatólogo; y (iv)~interfaz web multilingüe (catalán, español,
inglés) compatible con flujos de teledermatología.

La arquitectura resultante (sección~\ref{sec:dermapixel-arq}) implementa
una separación cliente-servidor con bus de mensajería intermedio
sobre AMQP. El cliente expone una API REST sobre FastAPI que
encapsula la lógica de persistencia y la sincronización
asíncrona; el servidor de inferencia opera sobre la DGX Spark con
todos los modelos cargados en VRAM unificada y consume mensajes
del bus.

\section{Arquitectura}
\label{sec:dermapixel-arq}

\subsection{Topología cliente-servidor}
\label{sec:dermapixel-topologia}

El sistema se distribuye en dos nodos físicos sobre una red local
(LAN). El primer nodo, equipo Windows del usuario clínico,
aloja la API REST, la base de datos, el
proxy inverso y la interfaz web. El segundo nodo, servidor GPU
on-prem NVIDIA DGX Spark con chip Grace Hopper GB10, aloja el
servidor de inferencia con los seis modelos cargados, el
servidor del modelo de lenguaje multimodal de gran tamaño y los
modelos de razonamiento auxiliares servidos vía Ollama. La
comunicación entre ambos nodos se realiza exclusivamente vía AMQP
sobre RabbitMQ, sin tráfico HTTP directo, lo que aísla
operativamente el servidor de inferencia tras el firewall del
proveedor de cómputo.

\subsection{Componentes del backend}
\label{sec:dermapixel-backend}

El backend (FastAPI $0.115$ asíncrono) se despliega como
contenedores Docker Compose con cinco servicios. La
tabla~\ref{tab:dermapixel-services} sintetiza puertos y
responsabilidades.

\begin{table}[H]
\centering
\footnotesize
\setlength{\tabcolsep}{4pt}
\renewcommand{\arraystretch}{0.95}
\caption{Servicios del backend DermApIxel.}
\label{tab:dermapixel-services}
\begin{tabular}{llrl}
\hline
Servicio                & Imagen               & Puerto host & Función \\
\hline
\texttt{derm-svc}       & FastAPI $0{,}115$ + Uvicorn  & $9030$ & API REST principal \\
\texttt{derm-mysql}     & MySQL $8.0$          & (interno) & Persistencia estructurada \\
\texttt{derm-rabbitmq}  & RabbitMQ $3.13$      & $5673 / 15675$ & Bus AMQP y panel de gestión \\
\texttt{derm-nginx}     & Nginx $1.25$         & $8091$ & Proxy inverso y archivos estáticos \\
\texttt{result-consumer}& Python asyncio       & ---          & Worker asíncrono que recibe resultados \\
\hline
\end{tabular}
\end{table}

\subsection{Persistencia: esquema de base de datos}
\label{sec:dermapixel-db}

La base de datos contiene doce tablas relacionales que documentan
cada estudio dermatológico, sus resultados parciales por módulo y
las validaciones posteriores. La tabla raíz \texttt{derm\_studies}
contiene el identificador único, el estado del estudio
($\texttt{pending}$, $\texttt{processing}$, $\texttt{done}$,
$\texttt{error}$), los \emph{timings} por etapa y la imagen
original almacenada como \texttt{LONGBLOB}. Las diez tablas
restantes contienen los resultados de cada módulo y se enlazan a
\texttt{derm\_studies} mediante clave foránea con borrado en
cascada. La separación por tabla, en lugar de un único documento
JSON, facilita las consultas analíticas posteriores (estadísticas
por clase, validación selectiva, exportación a CSV para
reentrenamiento) y mantiene la integridad referencial bajo cargas
concurrentes.

\subsection{Bus de mensajería AMQP}
\label{sec:dermapixel-amqp}

La comunicación entre backend e inferencia se realiza mediante dos
colas RabbitMQ con semántica distinta. La cola
\texttt{derm.inference} es \emph{work queue}: el backend publica
una tarea por estudio (imagen + identificador único) y el servidor
de inferencia la consume con confirmación explícita
(\emph{manual ack}) tras producir el resultado. La cola
\texttt{derm.results} es \emph{publish-subscribe} (configurada
como \emph{topic exchange}): el servidor de inferencia publica un
mensaje por estudio completado, y el worker \texttt{result-consumer}
del backend lo consume para actualizar la base de datos y notificar
al frontend. El módulo M6 (clasificación zero-shot) utiliza una
tercera vía RPC sobre AMQP para invocaciones síncronas de baja
latencia ($\sim 20$\,ms), apropiada para consultas
interactivas del clínico.

La elección de AMQP sobre HTTP directo aporta tres ventajas
operativas: tolerancia a fallos del servidor de inferencia (los
mensajes quedan en cola), desacoplamiento temporal entre cliente y
servidor (no se requiere conexión persistente), y trazabilidad
auditable de cada mensaje a través del panel de gestión de
RabbitMQ.

\section{Módulos de inteligencia artificial}
\label{sec:dermapixel-modulos}

El servidor de inferencia tiene seis módulos cargados
simultáneamente sobre la VRAM unificada de la DGX Spark. La
tabla~\ref{tab:dermapixel-modulos} sintetiza su correspondencia
con los modelos evaluados en el capítulo~\ref{cap:resultados}.

\begin{table}[H]
\centering
\footnotesize
\setlength{\tabcolsep}{3pt}
\renewcommand{\arraystretch}{0.95}
\caption{Módulos del prototipo DermApIxel.}
\label{tab:dermapixel-modulos}
\begin{tabular}{llrl}
\hline
Módulo & Modelo subyacente                       & Latencia & Función \\
\hline
M1     & PanDerm Large FT (HAM10000, TTA)         & $\sim 315$\,ms & Clasificación cerrada $7$ clases \\
M2     & SAM2.1-Large (decoder FT, ISIC2018)      & ---            & Segmentación binaria de lesión \\
M3     & SAE $1\,024 \to 16\,384$ + SkinCon       & ---            & $32$ conceptos clínicos por imagen \\
M4     & FAISS sobre Derm1M ($421\,327$ vec.)     & $\sim 150$\,ms & Recuperación visual densa \emph{top-5} \\
M5     & MedGemma~27B / GPT-4o / GPT-5 (9 prov.)  & $6$--$300$\,s  & Razonamiento clínico estructurado \\
M6     & DermLIP v2 (PubMedBERT)                  & $\sim 20$\,ms  & Clasificación abierta zero-shot \\
\hline
\end{tabular}
\end{table}

El módulo M7 (clasificador unificado L1/L2/L3 sobre el dataset
armonizado de $43$ clases L3 y M8 (SigLIP-Large SO400M
con sondeo lineal sobre HAM10000) están integrados en el pipeline
de evaluación pero no se exponen como pestañas independientes de
la interfaz clínica.

\subsection{M1 -- Clasificador cerrado}
\label{sec:dermapixel-m1}

M1 es PanDerm Large \emph{fine-tuned} sobre HAM10000 con la receta
descrita en sección~\ref{sec:ft}. La inferencia emplea \emph{test-time
augmentation} con cinco augmentaciones deterministas, lo que
incrementa la latencia respecto a la inferencia simple
(\emph{recall} de melanoma sobre HAM10000:
$+4{,}3$\,pp con TTA, según se documenta en la
sección~\ref{sec:ensemble-results}). La salida es un vector de siete
probabilidades softmax sobre las clases HAM10000, persistido en la
tabla \texttt{derm\_classifications}.

\subsection{M2 -- Segmentación binaria}
\label{sec:dermapixel-m2}

M2 es SAM2.1-Large adaptado mediante \emph{fine-tuning} de su
decodificador conforme a la
configuración descrita en sección~\ref{sec:segmentacion}. La inferencia
recibe la imagen y, en el modo interactivo, un \emph{prompt}
espacial (punto o caja delimitadora). La salida es una máscara
binaria almacenada como \texttt{LONGBLOB} en la tabla
\texttt{derm\_segmentations}, junto con métricas derivadas
(asimetría, área estimada en píxeles, perímetro). La interfaz
superpone la máscara sobre la imagen original como overlay
semitransparente.

\subsection{M3 -- Conceptos clínicos}
\label{sec:dermapixel-m3}

M3 es el módulo de interpretabilidad conceptual basado en Sparse
Autoencoders. El procedimiento es el siguiente: la imagen se
codifica mediante PanDerm Large, el \emph{embedding} de
$1\,024$ dimensiones se proyecta al espacio disperso de
$16\,384$ \emph{features} mediante el SAE entrenado sobre
$50\,454$ imágenes de siete datasets (sección~\ref{sec:disc-equidad}
para más detalle), y las activaciones se mapean a los $32$
conceptos clínicos de SkinCon~\cite{skincon} mediante regresión
logística por concepto. La salida es un vector de $32$
activaciones interpretables (\emph{pedunculado}, \emph{ulcerado},
\emph{eritematoso}, \emph{escamoso}, etc.) que se muestra al
clínico como apoyo visual al diagnóstico.

\subsection{M4 -- Recuperación visual densa}
\label{sec:dermapixel-m4}

M4 indexa los $421\,327$ \emph{embeddings} DermLIP de las
imágenes de Derm1M mediante FAISS~\cite{faiss2019} con índice HNSW
(\emph{Hierarchical Navigable Small Worlds}). La consulta produce
las cinco imágenes más similares en el espacio de
\emph{embeddings}, junto con su diagnóstico asignado y la distancia
coseno. La latencia media es de $\sim 150$ milisegundos por
consulta sobre el equipo descrito en sección~\ref{sec:infraestructura}.
La salida se persiste en \texttt{derm\_similar\_cases} y se
presenta como galería navegable de casos comparables.

La distinción operativa entre M4 y un sistema RAG completo es
relevante: M4 devuelve los casos análogos como referencia para
inspección humana, sin alimentar un modelo generativo posterior
con esos casos como contexto. La integración del retrieval con
generación se materializa parcialmente en M5
(sección~\ref{sec:dermapixel-m5}), donde los hallazgos de M1--M4 se
inyectan como contexto del prompt clínico.

\subsection{M5 -- Razonamiento clínico estructurado}
\label{sec:dermapixel-m5}

M5 genera una explicación textual en español que integra los
hallazgos de M1--M4 (y opcionalmente M6) en un razonamiento
clínico estructurado. La arquitectura soporta nueve proveedores
intercambiables mediante el parámetro \texttt{provider} del
endpoint correspondiente, cuyo perfil de coste y latencia se
sintetiza en la tabla~\ref{tab:dermapixel-llm}.

\begin{table}[H]
\centering
\footnotesize
\setlength{\tabcolsep}{3pt}
\renewcommand{\arraystretch}{0.95}
\caption{Proveedores LLM soportados por M5. Latencias medidas
sobre prompt clínico real ($\sim 1\,500$ tokens entrada,
$\sim 500$ tokens salida).}
\label{tab:dermapixel-llm}
\begin{tabular}{llrll}
\hline
\texttt{provider}        & Modelo                   & Latencia      & Coste/llamada & Infraestructura \\
\hline
\texttt{gpt-4o-mini}     & GPT-4o-mini               & $\sim 6$\,s  & $\sim$\$0{,}001 & API OpenAI \\
\texttt{gpt-4o}          & GPT-4o                    & $\sim 10$\,s & $\sim$\$0{,}01  & API OpenAI \\
\texttt{gpt-5-mini}      & GPT-5 mini                & $\sim 15$\,s & $\sim$\$0{,}001 & API OpenAI \\
\texttt{gpt-5-nano}      & GPT-5 nano                & $\sim 25$\,s & $\sim$\$0{,}0003 & API OpenAI \\
\texttt{gpt-5}           & GPT-5                     & $\sim 30$\,s & $\sim$\$0{,}01  & API OpenAI \\
\texttt{medgemma-4b}     & MedGemma~4B Vision        & $\sim 40$\,s & $0$\,EUR (local) & Ollama, $\sim 4$\,GB VRAM \\
\texttt{medgemma-vision} & MedGemma~27B Vision (BF16)& $3$--$5$\,min& $0$\,EUR (local) & DGX Spark, $54{,}9$\,GB VRAM \\
\hline
\end{tabular}
\end{table}

El \texttt{provider} por defecto para la interfaz interactiva es
\texttt{gpt-4o-mini} por su compromiso latencia-calidad. Los
proveedores locales (\texttt{medgemma-4b} y
\texttt{medgemma-vision}) se reservan para escenarios de
cumplimiento estricto sobre datos sanitarios donde la
exfiltración hacia APIs externas no es aceptable. La latencia de
los proveedores locales es superior a la de los proveedores
comerciales pese al tamaño nominalmente reducido, comportamiento
documentado en sección~\ref{sec:disc-llm} y atribuible a la diferencia
de optimización entre infraestructuras dedicadas y hardware
genérico.

Los nueve proveedores comparten el mismo \emph{system prompt} de
$\geq 300$ palabras estructurado en cinco secciones obligatorias
(morfología, dermatoscopia, criterios ABCDE con valores concretos
por letra, integración con resultados de M1--M4, juicio clínico),
y el mismo contexto enriquecido generado por la función
\texttt{\_build\_context()} que inyecta las cifras del pipeline al
\emph{prompt}. La salida es un JSON con tres campos:
\texttt{reasoning} (texto estructurado), \texttt{differential\_diagnosis}
(tres a cinco objetos con diagnóstico, nivel de confianza y
justificación) y \texttt{recommendation} (urgencia, acción
concreta, signos de alarma).

\subsection{M6 -- Clasificación zero-shot abierta}
\label{sec:dermapixel-m6}

M6 implementa la clasificación zero-shot multimodal descrita en
sección~\ref{sec:zs} usando DermLIP v2 (ViT-B/16 + PubMedBERT) con los
\emph{prompts} A3 derivados de Derm1M. A diferencia de M1
(clasificador cerrado sobre siete clases HAM10000), M6 acepta una
lista arbitraria de descripciones textuales y devuelve la similitud
coseno softmax sobre cada una. La latencia es de
$\sim 20$ milisegundos en \emph{warm} sobre los embeddings
precomputados.

El módulo expone cuatro \emph{presets} clínicos predefinidos
(tabla~\ref{tab:dermapixel-presets}) que cubren el espectro habitual
de las consultas dermatológicas, con un total de $31$ clases
únicas. La opción por defecto en la interfaz es el preset más
amplio (\texttt{all\_dermatology}), elegido para mitigar el sesgo
de selección documentado en la literatura sobre CLIP cuando se
restringe artificialmente el espacio de clases.

\begin{table}[H]
\centering
\footnotesize
\setlength{\tabcolsep}{4pt}
\renewcommand{\arraystretch}{0.95}
\caption{Presets clínicos del módulo M6.}
\label{tab:dermapixel-presets}
\begin{tabular}{lrl}
\hline
\texttt{preset}            & Clases & Uso clínico \\
\hline
\texttt{all\_dermatology}  & $31$   & Default: el clínico no precategoriza \\
\texttt{common\_dermatoses}& $8$    & Dirigido a sospecha de patología inflamatoria \\
\texttt{infections}        & $9$    & Dirigido a sospecha de infección \\
\texttt{tumors\_extended}  & $14$   & Dirigido a sospecha de tumor o lesión pigmentada \\
\hline
\end{tabular}
\end{table}

La nomenclatura y los recuentos de la tabla~\ref{tab:dermapixel-presets}
corresponden a la configuración inicial del prototipo; el catálogo de
\emph{presets} se amplió con posterioridad ---con selecciones por
subcategoría L2 y un ajuste de las claves de los presets
etiológicos--- sin alterar el diseño descrito.

Un hallazgo operativo documentado durante el despliegue
(sección~\ref{sec:disc-zs}) es la sensibilidad extrema de DermLIP v2 a
la formulación del \emph{prompt}: el uso de etiquetas
desnudas (e.g.\ \texttt{["melanoma", "nevus", "eczema"]}) produce
predicciones efectivamente aleatorias, mientras que el envoltorio
clínico estandarizado (e.g.~\emph{``dermoscopy image of melanoma,
malignant skin cancer with irregular borders and asymmetry''}) lo
restaura al rendimiento esperado. La curva de degradación
observada sobre HAM10000 sitúa el punto óptimo en
$8$--$15$ palabras por \emph{prompt} con un rasgo clínico
añadido (AUROC $0{,}854$); fuera de este rango el rendimiento
decae. El prototipo implementa \emph{wrapping} automático en tres
capas: presets predefinidos con descripciones ya envueltas,
composable \texttt{useZeroShot.js} que aplica la plantilla
\emph{``dermoscopy image of \{term\}''} cuando la entrada del
usuario contiene cinco o menos palabras, y conmutador en la
interfaz que permite al dermatólogo desactivar el envoltorio
cuando escribe descripciones completas propias.

\section{Frontend e interfaz clínica}
\label{sec:dermapixel-frontend}

El frontend está implementado en Vue~$3.5$ con Tailwind
CSS~$4.2$ y Vite~$7$, en tres idiomas (catalán, español,
inglés) y con tema claro/oscuro conmutable. La vista principal
\texttt{AnalyzeView} expone los resultados del pipeline en seis
pestañas: \texttt{CLASS} (clasificación cerrada M1),
\texttt{SEG} (segmentación M2 con overlay), \texttt{CONC}
(conceptos M3), \texttt{RAG} (casos similares M4),
\texttt{REASON} (razonamiento M5) y \texttt{Z-SHOT}
(clasificación abierta M6). Vistas complementarias
(\texttt{HistoryView}, \texttt{ValidationView},
\texttt{StatsView}) cubren la consulta del histórico paginado, la
validación posterior por el dermatólogo y las estadísticas
agregadas del sistema.

La interacción con el visor de lesiones
(\texttt{LesionViewer}) permite zoom real sobre la imagen,
superposición de la máscara con opacidad ajustable, y deslizador
de umbral dinámico que recalcula la predicción de M1 sobre la
marcha sin nuevas inferencias en el servidor. El módulo
\texttt{ValidationPanel} expone tres opciones de validación
(\emph{correcto}, \emph{parcial}, \emph{incorrecto}) y permite el
dibujo manual de \emph{bounding boxes} alternativas como mecanismo
de corrección, con persistencia en la tabla
\texttt{derm\_manual\_annotations} para realimentar el
reentrenamiento posterior.

\section{Rendimiento operativo}
\label{sec:dermapixel-rendimiento}

La tabla~\ref{tab:dermapixel-rendimiento} sintetiza los tiempos
medios medidos sobre el sistema desplegado.

\begin{table}[H]
\centering
\footnotesize
\setlength{\tabcolsep}{4pt}
\renewcommand{\arraystretch}{0.95}
\caption{Rendimiento operativo del prototipo DermApIxel (medidas
del 2026-04-05).}
\label{tab:dermapixel-rendimiento}
\begin{tabular}{lrl}
\hline
Métrica                                          & Valor               & Notas \\
\hline
Pipeline M1--M4 \emph{end-to-end}                 & $\sim 315$\,ms (warm) & Tras carga inicial \\
Pipeline M1--M4 \emph{cold}                       & $\sim 450$\,ms        & Primera inferencia tras reinicio \\
M6 Zero-shot ($3$ clases) \emph{warm}             & $\sim 20$\,ms server  & $\sim 25$\,ms incluyendo red \\
M6 Zero-shot ($3$ clases) \emph{cold}             & $\sim 25{,}5$\,ms     & $\sim 46{,}9$\,ms incluyendo red \\
M6 \emph{warmup} kernel PubMedBERT ($7$ clases)   & $614$\,ms             & Una vez tras reinicio \\
VRAM total app M1--M4 + M6                        & $4{,}40$\,GB         & Sin cuantización ni \emph{offloading} \\
VRAM MedGemma~27B Vision                          & $54{,}9$\,GB         & Coexiste con app principal \\
VRAM total ocupada                                & $\sim 59$\,GB        & Sobre $128$\,GB disponibles \\
Speedup M6 frente a M1--M4                        & $\sim 30\times$       & Por arquitectura ligera y \emph{embeddings} cacheados \\
\hline
\end{tabular}
\end{table}

La latencia agregada de un análisis completo
(M1--M4 + M5 con \texttt{gpt-4o-mini}) sobre una imagen de
prueba es del orden de $6$--$8$ segundos, dominada por el
componente generativo del razonamiento clínico. La latencia
percibida por el usuario es inferior gracias al patrón asíncrono:
los resultados de M1--M4 aparecen en la interfaz antes que el
razonamiento de M5, lo que permite al clínico revisar el
diagnóstico cerrado mientras el LLM produce la justificación
textual.

\section{Estado de despliegue}
\label{sec:dermapixel-estado}

A fecha de cierre del documento, el sistema se encuentra en
producción \emph{end-to-end} sobre el entorno de desarrollo
(Windows local + DGX Spark sobre red LAN), con seis módulos
operativos (M1--M6) y validación funcional sobre casos de prueba
controlados. El sistema no se ha desplegado sobre la
infraestructura sanitaria del HUSLL y no ha sido validado sobre
cohorte prospectiva real, lo que se identifica como línea de
continuación inmediata en sección~\ref{sec:linea-prospectivo}. Los
elementos identificados como condiciones operativas para el
despliegue hospitalario son: integración con el sistema PACS
mediante DICOM C-STORE, autenticación con Keycloak OIDC, HTTPS
con certificado institucional, servicio \texttt{systemd} para
arranque automático del servidor de inferencia, y trazabilidad
auditable conforme al Esquema Nacional de Seguridad (ENS) con
logs encadenados con SHA-256. Ninguno de estos elementos forma
parte del alcance del Trabajo de Fin de Grado declarado en
sección~\ref{sec:alcance}.

\section{Discusión: rol del prototipo en el trabajo}
\label{sec:dermapixel-discusion}

La construcción del prototipo no constituye objeto experimental
del trabajo. Las cifras de rendimiento sobre las que descansa la
defensa empírica del documento se obtienen directamente de los
modelos descritos en el capítulo~\ref{cap:metodos} mediante
protocolos estándar (sección~\ref{sec:protocolos}). El prototipo
materializa la integración técnica entre esos modelos, valida la
viabilidad operativa de su despliegue conjunto sobre hardware
accesible, y constituye el entorno previsto para la validación
clínica prospectiva descrita en sección~\ref{sec:linea-prospectivo}.

Tres aportaciones técnicas merecen mención específica. Primero,
la verificación de que seis modelos heterogéneos
(PanDerm Large FT, SAM2.1-Large decoder FT, SAE, FAISS sobre Derm1M,
DermLIP v2, MedGemma~27B) coexisten en la VRAM unificada de la
DGX Spark con consumo agregado de $\sim 59$\,GB sobre los
$128$\,GB disponibles, sin cuantización ni \emph{offloading}.
Segundo, la cuantificación del coste operativo por proveedor LLM
en M5, que sostiene la afirmación de sección~\ref{sec:disc-llm} sobre la
brecha entre encoders específicos y LLMs multimodales generalistas
en términos económicos además de en términos de rendimiento.
Tercero, la documentación operativa del hallazgo sobre
sensibilidad al \emph{prompt} en DermLIP v2
(sección~\ref{sec:disc-zs}), cuya mitigación mediante \emph{wrapping}
automático en tres capas (presets, composable, conmutador UI)
constituye una práctica de ingeniería transferible a otros
despliegues de modelos vision-lenguaje contrastivos en producción
clínica.

El despliegue hospitalario completo, la validación prospectiva
sobre cohorte real y la certificación CE como producto sanitario
quedan fuera del alcance declarado y se identifican como líneas
de continuación inmediatas en la sección~\ref{sec:lineas-abiertas} del
capítulo~\ref{cap:conclusiones}.

\section{Evolución a M1--M11: módulos castellanos del prototipo}
\label{sec:dermapixel-m11}

En los meses posteriores a la integración M1--M6 (estado
documentado en las secciones previas con fecha de cierre
2026-04-05), el prototipo incorporó tres módulos castellanos
adicionales (M9, M10, M4-bis) más un ensemble ponderado (M11)
sobre la ontología jerárquica de la Dra.~R.~Taberner descrita en
el Anexo~\ref{cap:anexoc}. La motivación se alinea con la
limitación operativa identificada en
sección~\ref{sec:dermapixel-discusion}: el clasificador cerrado
M1 sobre HAM10000 no cubre la \emph{fingerprint} clínica del
archivo de la Dra.~R.~Taberner, basada en un vocabulario de
$4$ categorías L1, $43$ subcategorías L2 y $367$ subcategorías
L3, mayoritariamente disjunto del de HAM10000. Los nombres M1b,
M7, M9, M10, M4-bis y M11 se conservan tal como aparecen en
\texttt{pipeline.py} para preservar la trazabilidad
código-documento.

\subsection{M7 -- Clasificador unificado jerárquico
\emph{merged43}+TTA}
\label{sec:dermapixel-m7}

M7 es PanDerm Large \emph{fine-tuned} sobre las $43$ clases L3 de
la ontología unificada multi-dataset, entrenado con
\emph{weighted sampling} sobre los siete datasets viables del
\emph{mapping} ontológico. La inferencia aplica \emph{Test-Time
Augmentation} (TTA) con cinco augmentaciones (original, simetría
horizontal, simetría vertical, rotación $90^\circ$ y rotación
$270^\circ$) promediando probabilidades \emph{softmax}, no
logits, por estabilidad numérica. La salida se proyecta
jerárquicamente a L1 y L2 garantizando consistencia top-down.
Sobre el conjunto de test fijo de DermapixelAI~1.0 alcanza
L1 \emph{accuracy} $0{,}947$, L2 \emph{accuracy} $0{,}819$,
L3 \emph{accuracy} $0{,}797$ y BAcc L3 $0{,}818$
($+4{,}72$\,pp BAcc con TTA respecto a la inferencia simple).
La latencia \emph{warm} es de $\sim 80$\,ms. Tab \texttt{UNIF}
del frontend, con vista jerárquica L1/L2/L3 y \emph{badge}
melanoma.

\subsection{M9 -- Dermapixel R0 (cabeza L2), clasificador L2
castellano supervisado}
\label{sec:dermapixel-m9}

M9 es la implementación operativa de la cabeza supervisada L2 de
Dermapixel R0
descrita en sección~\ref{sec:dermapixel-spanderm-v0}. Cabeza
\emph{fully-connected} $1\,024 \to 38$ sobre PanDerm Large con
LoRA $r = 16$ insertado en los dos últimos bloques transformer
(\texttt{attn.qkv}, \texttt{attn.proj}, \texttt{mlp.fc1},
\texttt{mlp.fc2}), entrenada sobre las $874$ imágenes \emph{train}
del manifiesto filtrado de DermapixelAI~1.0 ($1\,062$ imágenes;
Anexo~\ref{cap:anexoc}). Las $38$ clases L2
efectivas resultan de consolidar las dos variantes ortográficas
L2 \emph{queratinización} detectadas en el análisis exploratorio
del dataset. El \emph{checkpoint} \emph{slim} conserva únicamente
los pesos LoRA y la cabeza FC ($\sim 2{,}3$\,MB); el encoder
PanDerm Large permanece congelado y se reutiliza del módulo M1.
La inferencia aplica TTA con cinco augmentaciones y latencia
$\sim 91$\,ms. Tab \texttt{SPAN} del frontend. La diferencia
respecto a la rama contrastiva de Dermapixel R0 descrita en
sección~\ref{sec:spanderm-clip} es operacional: aquí la cabeza
es supervisada con clases L2 fijas; allí es \emph{zero-shot
open-class} por alineamiento texto-imagen.

\subsection{M10 -- \emph{Seven-Point Checklist} multitarea}
\label{sec:dermapixel-m10}

M10 implementa el protocolo multitarea de
sección~\ref{sec:sae-d7pt-e4} sobre el dataset
Derm7pt~\cite{kawahara2019}: \texttt{MultitaskDerm7} con siete
cabezas conceptuales (una por criterio dermatoscópico del
\emph{Seven-Point Checklist}: \emph{pigment network}, \emph{streaks},
\emph{pigmentation}, \emph{regression structures},
\emph{dots and globules}, \emph{blue whitish veil},
\emph{vascular structures}) más una cabeza binaria
melanoma/no-melanoma, montadas sobre el mismo encoder PanDerm
Large $+$ LoRA $r = 16$. La inferencia aplica TTA promediando
\emph{softmax} de las ocho cabezas y latencia $\sim 90$\,ms.
Tab \texttt{7-POI} del frontend, con los siete criterios
traducidos al castellano y el riesgo melanoma global. El módulo
explota explícitamente las anotaciones conceptuales por imagen
de Derm7pt, no aprovechadas en M1--M6 ni en otros experimentos
del trabajo.

\subsection{M4-bis -- RAG castellano sobre DermapixelAI 1.0}
\label{sec:dermapixel-m4bis}

M4-bis es un índice FAISS~\cite{faiss2019} \texttt{IndexFlatIP}
sobre $874$ \emph{embeddings} PanDerm Large de $1\,024$
dimensiones L2-normalizados, correspondientes al \emph{split}
\emph{train}$+$\emph{val} de DermapixelAI~1.0 (Anexo~\ref{cap:anexoc}).
Para una imagen de consulta el módulo devuelve los $k$ vecinos
más similares por coseno con sus metadatos clínicos castellanos
(\texttt{ontology\_l1}, \texttt{ontology\_l2}, \texttt{ontology\_l3},
\texttt{case\_id}, \texttt{case\_title}, \texttt{case\_text\_preview},
\texttt{rosa\_verified}, \texttt{year}) y el voto mayoritario por
nivel jerárquico. M4-bis complementa al módulo M4
(sección~\ref{sec:dermapixel-m4}) sobre Derm1M en inglés: donde
M4 recupera patrones visuales sobre un atlas masivo no etiquetado
en castellano, M4-bis devuelve casos del propio archivo de la
Dra.~R.~Taberner con su clasificación L1/L2/L3 verificada y el
extracto del texto clínico original. La latencia \emph{warm} es
de $\sim 0{,}5$\,ms (FAISS sobre $874$ vectores). Tab
\texttt{RAG-I} del frontend con barras de similitud y voto
mayoritario por nivel.

\subsection{M11 -- Ensemble ponderado por AUROC y coherencia
jerárquica}
\label{sec:dermapixel-m11sub}

M11 es la combinación ponderada (\emph{weighted}) de M1, M7, M9
y M4-bis por niveles L1/L2/L3. Los pesos se derivaron en dos
pasos: en primer lugar, lectura directa de las AUROC reportadas
en el capítulo~\ref{cap:resultados} (M7 \emph{merged43}+TTA es
el clasificador más fiable global, con AUROC L1 $= 0{,}873$,
L2 $= 0{,}860$, L3 $= 0{,}813$); en segundo lugar, ajuste empírico
por coherencia jerárquica top-down sobre casos clínicos
representativos. La iteración estable es la versión \texttt{v1.7},
cuyos pesos se reportan en la tabla~\ref{tab:m11-weights}.

\begin{table}[H]
\centering
\footnotesize
\setlength{\tabcolsep}{4pt}
\renewcommand{\arraystretch}{0.95}
\caption{Pesos del ensemble M11 \texttt{v1.7} por nivel jerárquico.
La entrada \emph{(no aplica)} indica que el módulo correspondiente
no emite predicción al nivel ontológico de la fila.}
\label{tab:m11-weights}
\begin{tabular}{lrrrr}
\hline
Nivel & M1 (HAM FT) & M7 (\emph{merged43}+TTA) & M9 (Dermapixel R0) & M4-bis (RAG ES) \\
\hline
L1   & $1{,}0$           & $1{,}0$           & $0{,}5$           & $0{,}3$           \\
L2   & \emph{(no aplica)}& $1{,}5$           & $1{,}0$           & $0{,}5$           \\
L3   & \emph{(no aplica)}& $1{,}5$           & \emph{(no aplica)}& $0{,}5$           \\
\hline
\end{tabular}
\end{table}

El \emph{mapping} castellano-inglés exhaustivo ($4$ L1
$+$ $26$ L2 $+$ $33$ L3) colapsa las duplicaciones que aparecen
cuando M7 emite la ontología en inglés (\emph{merged43}) mientras
que M9 y M4-bis lo hacen en castellano. Sin este \emph{mapping},
el ensemble contaría dos veces la misma clase con cadenas
distintas. La coherencia jerárquica queda garantizada: para una
imagen donde L1 indica \emph{Patología tumoral} con probabilidad
elevada y L3 indica \emph{Melanoma} con probabilidad relevante,
L2 debe indicar \emph{Tumores malignos} y no \emph{Neoplasias
melanocíticas benignas} (este último era el comportamiento del
ensemble cuando los pesos eran M7 $= 1{,}0$ y M9 $= 1{,}5$ en L2,
antes del ajuste \texttt{v1.7}). M11 se expone en el tab
\texttt{M11} del frontend como vista de consenso, posicionada en
primer lugar.

\subsection{Sistema de consenso de cinco módulos y \emph{warning}
de derivación urgente}
\label{sec:dermapixel-consenso}

El frontend incorpora un \emph{banner} superior con cinco
\emph{chips} que resumen la lectura de malignidad de cinco
módulos independientes (M1, M7, M9, SigLIP-LP y M10) más un sexto
\emph{chip} M4-bis con la ontología del vecino \emph{top-1}
castellano. Cada \emph{chip} muestra una etiqueta \emph{BEN}
/ \emph{MAL} / \emph{?} y un \emph{badge} global \emph{Consenso
parcial (X/5)} o \emph{Consenso fuerte}. Cuando $\geq 2$
módulos votan melanoma con probabilidad alta, el sistema emite
un \emph{warning} explícito (``X clasificadores detectan
melanoma -- derivación urgente -- derivar a especialista''). El
diseño replica el patrón clínico de segunda opinión entre
dermatólogos: no se confía en un único clasificador, se busca
convergencia entre criterios independientes (cabeza supervisada
HAM, cabeza supervisada L3 multidataset, cabeza supervisada L2
castellana, sondeo lineal SigLIP, multitarea Derm7pt). Un
\emph{caveat} metodológico explícito acompaña la formulación: el
voto actual es por mayoría sin ponderación clínica, lo que no
captura el caso en el que un módulo emite una señal dominante
(\emph{p.\,ej.} M10 con probabilidad $0{,}95$ de melanoma) frente
a otros con confianza baja. Una iteración futura debería
ponderar por confianza per-módulo o exponer un \emph{score} de
melanoma agregado calibrado con datos prospectivos.

\subsection{Estrategias de seguridad clínica: ensemble para detección de melanoma}
\label{sec:ensemble-results}

Más allá de la comparación de modelos individuales, resulta
relevante explorar si la complementariedad entre arquitecturas puede
explotarse para maximizar la sensibilidad diagnóstica. Esta cuestión
es especialmente importante en melanoma, donde el coste clínico de
un falso negativo es potencialmente muy superior al de un falso
positivo. Por ello se evalúa un \emph{ensemble} orientado
explícitamente a maximizar el \emph{recall} de melanoma sobre
HAM10000.

Se combinan tres clasificadores independientes sobre el
\emph{split} de test de HAM10000 ($N = 1\,232$ imágenes,
$70$ melanomas):
M1 (PanDerm Large \emph{fine-tuned} sobre HAM10000), SigLIP-Large
SO400M con sondeo lineal sobre los \emph{embeddings} de
$1\,152$ dimensiones, y M7 (clasificador unificado sobre el
\emph{merged43} del dataset armonizado, descrito en el
Anexo~\ref{cap:anexoh}).

\begin{table}[H]
\centering
\footnotesize
\setlength{\tabcolsep}{4pt}
\renewcommand{\arraystretch}{0.95}
\caption{Ensemble sobre HAM10000 ($N = 1\,232$, $70$ melanomas).
\emph{Recall} \emph{top-3}: la clase melanoma aparece entre las
tres primeras del ranking. FP\,mel: falsos positivos melanoma
(imágenes no-melanoma marcadas).}
\label{tab:ensemble}
\begin{tabular}{lrrrrrr}
\hline
Configuración                   & Acc           & BAcc          & AUROC         & Mel rec.\ top-1 & Mel rec.\ top-3 & FP\,mel \\
\hline
M1 (PanDerm FT)                 & $0{,}919$     & $0{,}813$     & $0{,}986$     & $0{,}714$ (50/70) & $0{,}971$ (68/70) & $36$ \\
SigLIP-L LP                     & $0{,}900$     & $0{,}776$     & $0{,}971$     & $0{,}614$ (43/70) & $0{,}986$ (69/70) & $31$ \\
Avg M1 + SigLIP                 & $0{,}917$     & $0{,}836$     & $0{,}986$     & $0{,}686$         & $0{,}986$         & $32$ \\
MaxMel M1 + SigLIP              & $0{,}909$     & $0{,}832$     & $0{,}985$     & $\mathbf{0{,}771}$ & $\mathbf{1{,}000}$ & $50$ \\
OR top-3 M1 + SigLIP            & ---           & ---           & ---           & ---               & $\mathbf{1{,}000}$ (70/70) & $885$ \\
\hline
\end{tabular}
\end{table}

La combinación OR \emph{top-3} de M1 y SigLIP alcanza un
\emph{recall} de melanoma del $100\%$ ($70$ de $70$), pero a costa
de $885$ falsos positivos sobre las $1\,232$ imágenes del test: la
sensibilidad máxima se obtiene marcando como sospechosa cerca del
$72\%$ del conjunto, lo que reduce drásticamente la precisión y
acota su utilidad a un escenario de \emph{cribado}, nunca de
decisión diagnóstica.

El principal hallazgo de esta sección, sin embargo, no es ese
$100\%$ ---una cifra trivialmente alcanzable ampliando el margen de
sospecha de cualquier clasificador--- sino la evidencia de
complementariedad entre una arquitectura especializada (M1, PanDerm
\emph{fine-tuned} \emph{in-domain}) y una generalista (SigLIP con
sondeo lineal): ambas recuperan melanomas distintos. Los dos
melanomas que M1 omite en \emph{top-3} (ISIC\_0024886.jpg e
ISIC\_0030360.jpg) sí los recupera SigLIP, y el único melanoma que
SigLIP no sitúa en \emph{top-3} sí lo recupera M1. Este resultado
indica que los errores de ambos modelos no están completamente
correlacionados ---condición necesaria para que una estrategia de
\emph{ensemble} aporte un beneficio clínico real, y no una mera
redundancia--- y sugiere que cada arquitectura explota
representaciones parcialmente diferentes de la lesión. M7, evaluado
sobre el subconjunto de $198$ imágenes con predicción simultánea de
los tres modelos ($12$ melanomas), no recupera melanomas adicionales
sobre el par, pero aporta cobertura sobre las $43$ clases
ontológicas del dataset armonizado.

Tres caveats acompañan la interpretación de este resultado: (i)~el
\emph{recall} del $100\%$ es \emph{top-3}, no \emph{top-1};
(ii)~el tamaño muestral de melanomas es $N = 70$, con intervalos
de confianza al $95\%$ amplios; (iii)~no existe validación
prospectiva sobre cohorte hospitalaria. El resultado se interpreta
como prueba de concepto para un escenario de cribado oportunista,
no como sustituto del juicio clínico.

\subsection{Migración a AMQP (2026-04-13)}
\label{sec:dermapixel-amqp-migracion}

Desde el 13 de abril de 2026 toda comunicación entre el backend
Docker (PC Windows) y el servidor de inferencia (DGX Spark)
circula por RabbitMQ (cola \texttt{derm.inference} para
\emph{requests}, \texttt{derm.results} para \emph{replies}).
Anteriormente coexistían HTTP directo y AMQP; tras la migración,
HTTP queda reservado a \texttt{/health} y depuración. La
configuración mantiene abierto el puerto $5673$ en dirección
backend $\to$ Spark. Los beneficios operativos son tres:
resiliencia (los mensajes persisten ante caída transitoria del
servidor de inferencia con reintento automático), trazabilidad
auditable de cada \emph{request} a través del panel de gestión
de RabbitMQ, y desacoplamiento temporal entre frontend y GPU. La
lección operacional documentada es que tras cualquier cambio
del backend Docker es obligatorio reconstruir sin caché
(\texttt{docker compose build --no-cache \&\& docker compose up
-d}) para que los nuevos \emph{consumers} AMQP se registren
correctamente.

\subsection{Evaluación end-to-end del prototipo extendido}
\label{sec:dermapixel-batch-test}

Para validar el comportamiento integrado de los módulos M1--M11
se ejecutó un \emph{batch test} sobre $454$ imágenes del
conjunto de test de HAM10000 no vistas previamente por el
prototipo. Los hallazgos cuantitativos justifican operacionalmente
los pesos del ensemble M11 \texttt{v1.7} y la presencia de
SigLIP-LP como red de seguridad. M7 supera a M1 sobre HAM10000
con BAcc $0{,}886$ frente a BAcc $0{,}743$, lo que justifica el
peso M7 $= 1{,}5$ en L2/L3 del ensemble a pesar de que M1 está
\emph{fine-tuned} específicamente sobre HAM. El \emph{recall} de
melanoma por módulo es de $40\%$ para M1, $60\%$ para M7 y
$80\%$ para el OR \emph{top-3} entre M1, M7 y SigLIP-LP, lo que
justifica directamente el patrón de consenso de cinco módulos
como red de seguridad clínica. El primer nivel jerárquico
(Tumoral / Inflamatoria / Infecciosa / Genodermatosis) alcanza
\emph{accuracy} en el intervalo $[0{,}80,\,1{,}00]$ sobre las
cuatro categorías, mientras que los errores \emph{fine-grained}
se concentran en L2 y L3, comportamiento coherente con la
naturaleza progresivamente más específica de la jerarquía
ontológica. TTA añade $+74$\,ms de latencia por imagen y
$+4{,}72$\,pp de BAcc L3, compromiso latencia-rendimiento
replicado además en M9 y M10. La evidencia operacional cierra
el ciclo entre las métricas del capítulo~\ref{cap:resultados} y
el comportamiento del sistema integrado: los módulos castellanos
M7--M11 mejoran las métricas críticas (\emph{recall} de melanoma,
BAcc sobre L2/L3) en el escenario clínico final sin sustituir
a M1--M6.
